\documentclass[a4paper,10pt]{article}

% Настройки русского языка и кодировок
\usepackage[T2A]{fontenc}
\usepackage[utf8]{inputenc}
\usepackage[russian]{babel}

% Математические пакеты
\usepackage{amsmath, amsfonts, amssymb, amsthm}
\usepackage{geometry}
\usepackage{titlesec}

% Настройка полей (компактные поля для конспекта)
\geometry{left=1cm, right=1cm, top=1.5cm, bottom=1.5cm}

% Настройка заголовков
\titleformat{\section}{\large\bfseries}{}{0em}{}
\titleformat{\subsection}{\normalsize\bfseries}{}{0em}{}

\begin{document}

\begin{center}
    \LARGE \textbf{Начально-краевая задача на отрезке. Метод Фурье}
\end{center}

\section{Часть 1: Практическая (Алгоритмы и факты)}

\subsection*{Постановка задачи}
Рассматривается начально-краевая задача (НКЗ) для функции $u(x, t) \in C^2(G) \cap C^1(\overline{G})$, где область $G = \{0 < x < l, \ t > 0\}$. Алгоритм един для волнового уравнения (ВУ) и уравнения теплопроводности (УТ). Заданы дифференциальное уравнение, начальные и граничные (краевые) условия:
\[
\begin{cases}
u_{tt} \ \text{(или } u_t\text{)} = a^2 u_{xx} + f(x,t) & (1) \quad 0 < x < l, \ t > 0 \\
u|_{t=0} = u_0(x), \quad \left[ u_t|_{t=0} = u_1(x) \text{ --- только для ВУ} \right] & (2) \quad 0 \le x \le l \\
(\alpha u - \beta u_x)|_{x=0} = \mu(t) & (3) \quad t \ge 0 \\
(\gamma u + \delta u_x)|_{x=l} = \nu(t) & (4) \quad t \ge 0
\end{cases}
\]
Предполагается, что константы $\alpha, \beta, \gamma, \delta \ge 0$, причем $|\alpha| + |\beta| > 0$ и $|\gamma| + |\delta| > 0$.

\subsection*{Алгоритм решения (Метод разделения переменных)}

\textbf{Шаг 0. Сведение задачи к однородным краевым условиям} \\
Решение ищется в виде суммы $u(x,t) = v(x,t) + w(x,t)$, где вспомогательная функция $w(x,t)$ подбирается так, чтобы она удовлетворяла неоднородным краевым условиям (3) и (4). Тогда новая функция $v(x,t)$ будет иметь нулевые краевые условия.
Выбор формы $w(x,t)$ зависит от типа краевых условий:
\begin{itemize}
    \item \textbf{Особый случай ($\alpha = 0, \gamma = 0$):} Заданы производные на обоих концах (задача Неймана). Функция ищется в виде полинома второй степени по $x$:
    \[ w(x,t) = x( \varphi(t) + x \psi(t) ) \]
    \item \textbf{Неособый случай (все остальные варианты):} Хотя бы на одном конце функция входит в условие с ненулевым коэффициентом. Функция ищется в виде линейной функции по $x$:
    \[ w(x,t) = \varphi(t) + x \psi(t) \]
\end{itemize}
Функции $\varphi(t)$ и $\psi(t)$ однозначно находятся путем подстановки выбранного вида $w(x,t)$ в условия (3) и (4) (см. обоснование в Части 2).

\vspace{0.2cm}
\textbf{Шаг 1. Формирование задачи для $v(x,t)$ и задача Штурма-Лиувилля} \\
После подстановки $u = v + w$ в исходную систему (1)-(4) получаем новую задачу для $v(x,t)$:
\[
\begin{cases}
v_{tt} \ \text{(или } v_t\text{)} = a^2 v_{xx} + g(x,t) \\
v|_{t=0} = v_0(x), \quad \left[ v_t|_{t=0} = v_1(x) \right] \\
(\alpha v - \beta v_x)|_{x=0} = 0, \quad (\gamma v + \delta v_x)|_{x=l} = 0
\end{cases}
\]
где $g(x,t) = f(x,t) + a^2 w_{xx} - w_{tt}$ (или $- w_t$ для УТ), $v_0(x) = u_0(x) - w(x,0)$, $v_1(x) = u_1(x) - w_t(x,0)$.

Для однородного уравнения $v_{tt} = a^2 v_{xx}$ с нулевыми краевыми условиями применяется метод разделения переменных $v(x,t) = T(t)X(x)$. Это приводит к \textbf{спектральной задаче Штурма-Лиувилля}:
\[
\begin{cases}
-X''(x) = \lambda X(x), \quad 0 < x < l \\
\alpha X(0) - \beta X'(0) = 0 \\
\gamma X(l) + \delta X'(l) = 0
\end{cases}
\]
\textbf{Свойства задачи Штурма-Лиувилля:}
\begin{enumerate}
    \item Множество собственных значений (СЗ) $\lambda_n$ счетно.
    \item Все СЗ неотрицательны: $\lambda_n \ge 0$ (в неособом случае $\lambda_n > 0$).
    \item Все СЗ простые (каждому $\lambda_n$ соответствует ровно одна собственная функция $X_n(x)$ с точностью до константы).
    \item $\lambda_n \to +\infty$ при $n \to \infty$. Упорядочиваются по возрастанию: $\lambda_1 < \lambda_2 < \lambda_3 < \dots$
    \item Система собственных функций $\{X_n(x)\}_{n=1}^\infty$ образует ортогональный базис в пространстве $L_2(0, l)$.
\end{enumerate}
Необходимо решить задачу Штурма-Лиувилля: найти все $\lambda_n$ и соответствующие им $X_n(x)$.

\vspace{0.2cm}
\textbf{Шаг 2. Разложение в ряды Фурье по базису $\{X_n(x)\}$} \\
Функции $v_0(x)$, $v_1(x)$ и $g(x,t)$ раскладываются в обобщенные ряды Фурье:
\[ v_0(x) = \sum_{n=1}^\infty a_n X_n(x), \quad v_1(x) = \sum_{n=1}^\infty b_n X_n(x), \quad g(x,t) = \sum_{n=1}^\infty c_n(t) X_n(x) \]
где коэффициенты вычисляются по формулам:
\[ a_n = \frac{1}{\|X_n\|^2} \int_0^l v_0(x) X_n(x) dx, \quad b_n = \frac{1}{\|X_n\|^2} \int_0^l v_1(x) X_n(x) dx, \quad c_n(t) = \frac{1}{\|X_n\|^2} \int_0^l g(x,t) X_n(x) dx \]
Квадрат нормы базисной функции вычисляется как $\|X_n\|^2 = \int_0^l X_n^2(x) dx$.

\vspace{0.2cm}
\textbf{Шаг 3. Решение обыкновенных дифференциальных уравнений (ОДУ)} \\
Решение $v(x,t)$ ищется в виде ряда $v(x,t) = \sum_{n=1}^\infty T_n(t) X_n(x)$. 
\begin{itemize}
    \item \textbf{Для волнового уравнения:} функция $T_n(t)$ удовлетворяет задаче Коши:
    \[ \begin{cases} T_n''(t) + a^2 \lambda_n T_n(t) = c_n(t) \\ T_n(0) = a_n, \quad T_n'(0) = b_n \end{cases} \]
    \item \textbf{Для уравнения теплопроводности:} функция $T_n(t)$ удовлетворяет задаче Коши:
    \[ \begin{cases} T_n'(t) + a^2 \lambda_n T_n(t) = c_n(t) \\ T_n(0) = a_n \end{cases} \]
\end{itemize}
\textit{Замечание:} Если исходное уравнение имело более общий вид $u_{tt} + b u_t + c u = a^2 u_{xx} + f(x,t)$, то уравнение для $T_n(t)$ модифицируется: $T_n''(t) + b T_n'(t) + (c + a^2 \lambda_n)T_n(t) = c_n(t)$.

\vspace{0.2cm}
\textbf{Шаг 4. Сборка итогового решения} \\
Решаются все ОДУ для $T_n(t)$, после чего записывается итоговый ответ:
\[ u(x,t) = \sum_{n=1}^\infty T_n(t) X_n(x) + w(x,t) \]

\newpage
\section{Часть 2: Теоретическая (Обоснования и доказательства)}

\subsection*{Почему метод разделения переменных работает и нет других решений?}
Метод ищет решения в виде разделенных переменных $T(t)X(x)$, что кажется ограничением (почему бы не искать в виде $e^{xt}$ и т.д.?). Строгое обоснование проводится в \textbf{спектральной теории операторов}.
Оператор $\frac{d^2}{dx^2}$ с заданными краевыми условиями является самосопряженным в гильбертовом пространстве $L_2(0,l)$. По теореме Стеклова, его собственные функции $\{X_n(x)\}$ образуют \textbf{полный ортогональный базис} в $L_2(0,l)$. 
Это означает, что абсолютно \textit{любая} достаточно гладкая функция (включая начальные условия $u_0(x), u_1(x)$ и правые части $f(x,t)$ при каждом фиксированном $t$) может быть однозначно разложена в сходящийся ряд по этому базису. Следовательно, решение, представленное в виде ряда Фурье $u(x,t) = \sum T_n(t)X_n(x)$, исчерпывает всё пространство возможных решений задачи Коши-Дирихле/Неймана. Никаких «потерянных» решений других видов (типа $e^{xt}$) не существует, так как они в любом случае проецируются на этот базис.

\subsection*{Обоснование Шага 0: Построение функции $w(x,t)$}

\textbf{Утверждение 1.} \textit{В неособом случае ($|\alpha| + |\gamma| > 0$) функцию $w(x,t)$, удовлетворяющую граничным условиям, можно найти в виде $w(x,t) = \varphi(t) + x \psi(t)$.}

\begin{proof}
Подставим предполагаемый вид $w(x,t)$ и ее производную $w_x(x,t) = \psi(t)$ в неоднородные граничные условия (3) и (4):
\[
\begin{cases}
\alpha (\varphi(t) + 0 \cdot \psi(t)) - \beta \psi(t) = \mu(t) \quad \text{(при } x=0) \\
\gamma (\varphi(t) + l \psi(t)) + \delta \psi(t) = \nu(t) \quad \text{(при } x=l)
\end{cases}
\]
Перепишем систему относительно неизвестных функций $\varphi(t)$ и $\psi(t)$ в матричном виде:
\[
\begin{pmatrix} 
\alpha & -\beta \\ 
\gamma & \gamma l + \delta 
\end{pmatrix} 
\begin{pmatrix} 
\varphi(t) \\ 
\psi(t) 
\end{pmatrix} 
= 
\begin{pmatrix} 
\mu(t) \\ 
\nu(t) 
\end{pmatrix}
\]
Вычислим определитель матрицы системы $\Delta$:
\[
\Delta = \alpha(\gamma l + \delta) - (-\beta\gamma) = \alpha\gamma l + \alpha\delta + \beta\gamma.
\]
Поскольку константы $\alpha, \beta, \gamma, \delta \ge 0$ и длина отрезка $l > 0$, проверим знак $\Delta$:
1. Если $\alpha > 0$, то либо $\gamma > 0$ (тогда $\alpha\gamma l > 0$), либо $\gamma = 0$. Если $\gamma = 0$, то по условию неособого случая $|\gamma| + |\delta| > 0 \implies \delta > 0$, тогда $\alpha\delta > 0$. Следовательно, $\Delta > 0$.
2. Если $\alpha = 0$, то по условию $|\alpha| + |\beta| > 0 \implies \beta > 0$. Поскольку рассматривается неособый случай, $\gamma$ не может быть равно нулю одновременно с $\alpha$, значит $\gamma > 0$. Тогда $\beta\gamma > 0$, и $\Delta > 0$.

Таким образом, в неособом случае определитель $\Delta$ всегда строго больше нуля. Система имеет единственное решение, и функции $\varphi(t)$ и $\psi(t)$ корректно определены.
\end{proof}

\textbf{Утверждение 2.} \textit{В особом случае ($\alpha = 0, \gamma = 0$) функцию $w(x,t)$ можно найти в виде $w(x,t) = x \varphi(t) + x^2 \psi(t)$.}

\begin{proof}
При $\alpha = 0, \gamma = 0$ краевые условия принимают вид:
\[
- \beta w_x(0,t) = \mu(t), \quad \delta w_x(l,t) = \nu(t).
\]
Отсюда получаем явные требования на производную на концах отрезка (заметим, что по условию невырожденности $\beta > 0$ и $\delta > 0$):
\[
w_x(0,t) = -\frac{\mu(t)}{\beta}, \quad w_x(l,t) = \frac{\nu(t)}{\delta}.
\]
Возьмем $w(x,t) = x \varphi(t) + x^2 \psi(t)$. Ее производная по пространственной координате равна $w_x(x,t) = \varphi(t) + 2x \psi(t)$.
Подставим $x=0$:
\[
w_x(0,t) = \varphi(t) = -\frac{\mu(t)}{\beta}.
\]
Функция $\varphi(t)$ найдена явно. Теперь подставим $x=l$:
\[
w_x(l,t) = \varphi(t) + 2l \psi(t) = \frac{\nu(t)}{\delta}.
\]
Поскольку $l > 0$, мы можем выразить $\psi(t)$:
\[
\psi(t) = \frac{1}{2l} \left( \frac{\nu(t)}{\delta} - \varphi(t) \right) = \frac{1}{2l} \left( \frac{\nu(t)}{\delta} + \frac{\mu(t)}{\beta} \right).
\]
Обе функции найдены, следовательно, указанный вид $w(x,t)$ применим для построения вспомогательной функции в особом случае.
\end{proof}

\subsection*{Обоснование Шага 1: Происхождение задачи Штурма-Лиувилля}

\begin{proof}
Рассмотрим однородное уравнение (например, волновое $v_{tt} = a^2 v_{xx}$) с однородными краевыми условиями. Метод Фурье заключается в поиске нетривиального решения в виде произведения двух функций, каждая из которых зависит только от одной переменной: $v(x,t) = T(t)X(x)$.
Подставим это представление в уравнение:
\[ T''(t)X(x) = a^2 T(t)X''(x). \]
Разделим обе части на $a^2 T(t)X(x)$ (предполагая, что оно не равно нулю):
\[ \frac{T''(t)}{a^2 T(t)} = \frac{X''(x)}{X(x)}. \]
Левая часть полученного равенства зависит только от $t$, а правая — только от $x$. Это тождество может выполняться при любых независимых $x$ и $t$ только в том случае, если обе части равны одной и той же константе. Обозначим эту константу (константу разделения) через $-\lambda$:
\[ \frac{X''(x)}{X(x)} = -\lambda \implies -X''(x) = \lambda X(x). \]
Теперь подставим $v(x,t) = T(t)X(x)$ в левое однородное краевое условие при $x=0$:
\[ \alpha v(0,t) - \beta v_x(0,t) = \alpha T(t)X(0) - \beta T(t)X'(0) = T(t) \big(\alpha X(0) - \beta X'(0)\big) = 0. \]
Так как мы ищем нетривиальное решение, $T(t) \not\equiv 0$, следовательно, от пространственной функции требуется выполнение условия $\alpha X(0) - \beta X'(0) = 0$.
Абсолютно аналогично из правого краевого условия ($x=l$) получаем $\gamma X(l) + \delta X'(l) = 0$.
Таким образом, для функции $X(x)$ естественным образом возникает классическая задача на собственные значения — задача Штурма-Лиувилля. (Для уравнения теплопроводности рассуждения аналогичны, лишь в левой части будет стоять $\frac{T'(t)}{a^2 T(t)}$).
\end{proof}

\subsection*{Доказательство свойств задачи Штурма-Лиувилля}

\textbf{Свойство 2 (Неотрицательность собственных значений).} \textit{Все собственные значения $\lambda \ge 0$.}
\begin{proof}
Пусть $X(x)$ --- собственная функция, соответствующая собственному значению $\lambda$. Умножим исходное уравнение $-X''(x) = \lambda X(x)$ на $X(x)$ и проинтегрируем по отрезку $[0, l]$:
\[ -\int_0^l X''(x)X(x) dx = \lambda \int_0^l X^2(x) dx. \]
К левой части применим формулу интегрирования по частям:
\[ -\left( X'(x)X(x) \big|_0^l - \int_0^l (X'(x))^2 dx \right) = \int_0^l (X'(x))^2 dx - X'(l)X(l) + X'(0)X(0). \]
Рассмотрим внеинтегральные члены, используя граничные условия:
1. При $x=l$: $\gamma X(l) + \delta X'(l) = 0$. Если $\delta > 0$, то $X'(l) = -\frac{\gamma}{\delta}X(l)$, и $-X'(l)X(l) = \frac{\gamma}{\delta}X^2(l) \ge 0$. Если $\delta = 0$, то $X(l) = 0$, и член равен $0$.
2. При $x=0$: $\alpha X(0) - \beta X'(0) = 0$. Если $\beta > 0$, то $X'(0) = \frac{\alpha}{\beta}X(0)$, и $X'(0)X(0) = \frac{\alpha}{\beta}X^2(0) \ge 0$. Если $\beta = 0$, то $X(0) = 0$, и член равен $0$.

Следовательно, оба внеинтегральных члена неотрицательны. Получаем:
\[ \lambda \int_0^l X^2(x) dx = \int_0^l (X'(x))^2 dx + \left[ \text{неотрицательные граничные слагаемые} \right] \ge 0. \]
Так как собственная функция $X(x) \not\equiv 0$, то $\int_0^l X^2(x) dx > 0$. Отсюда с необходимостью следует, что $\lambda \ge 0$. 
\textit{Замечание:} Равенство $\lambda = 0$ возможно только если интеграл от квадрата производной равен нулю ($X(x) = \text{const}$) и оба внеинтегральных члена равны нулю. Это реализуется только в особом случае (задача Неймана: $\alpha = 0, \gamma = 0$).
\end{proof}

\textbf{Свойство 3 (Простота собственных значений).} \textit{Каждому собственному значению соответствует ровно одна (с точностью до умножения на константу) собственная функция.}
\begin{proof}
Пусть одному и тому же значению $\lambda$ соответствуют две собственные функции $X_1(x)$ и $X_2(x)$. Рассмотрим их определитель Вронского: $W(x) = X_1(x)X_2'(x) - X_1'(x)X_2(x)$.
Вычислим его производную:
\[ W'(x) = X_1'(x)X_2'(x) + X_1(x)X_2''(x) - X_1''(x)X_2(x) - X_1'(x)X_2'(x) = X_1(x)(-\lambda X_2(x)) - (-\lambda X_1(x))X_2(x) = 0. \]
Следовательно, $W(x) \equiv \text{const}$. Вычислим константу при $x=0$.
Из граничного условия $\alpha X_i(0) - \beta X_i'(0) = 0$ следует, что векторы $(X_1(0), X_1'(0))$ и $(X_2(0), X_2'(0))$ ортогональны одному и тому же ненулевому вектору $(\alpha, -\beta)$. В двумерном пространстве это означает, что они коллинеарны, то есть их определитель равен нулю:
\[ W(0) = X_1(0)X_2'(0) - X_1'(0)X_2(0) = 0. \]
Раз $W(0) = 0$ и $W(x) \equiv \text{const}$, то $W(x) \equiv 0$ на всем отрезке $[0,l]$. Равенство нулю вронскиана решений уравнения второго порядка означает их линейную зависимость. Следовательно, $X_1(x) = C \cdot X_2(x)$.
\end{proof}

\textbf{Свойство 5 (Ортогональность базиса).} \textit{Собственные функции, соответствующие различным собственным значениям, ортогональны в $L_2(0,l)$.}
\begin{proof}
Пусть $\lambda_n \neq \lambda_m$ — два различных СЗ, а $X_n$ и $X_m$ — их собственные функции. Запишем уравнения:
\[ -X_n'' = \lambda_n X_n \]
\[ -X_m'' = \lambda_m X_m \]
Умножим первое уравнение на $X_m$, второе на $X_n$, вычтем одно из другого и проинтегрируем по $[0,l]$:
\[ (\lambda_n - \lambda_m) \int_0^l X_n(x)X_m(x) dx = \int_0^l \big( X_m''(x)X_n(x) - X_n''(x)X_m(x) \big) dx. \]
Подынтегральное выражение в правой части является точной производной: $(X_m' X_n - X_n' X_m)'$. Применяя формулу Ньютона-Лейбница, получаем:
\[ (\lambda_n - \lambda_m) \int_0^l X_n X_m dx = \big[ X_m'(x)X_n(x) - X_n'(x)X_m(x) \big]_0^l. \]
В силу граничных условий (аналогично доказательству простоты СЗ через коллинеарность), выражение в квадратных скобках (которое по сути является определителем Вронского для этих двух функций) обращается в ноль и на левом, и на правом концах.
Таким образом, правая часть равна нулю. Так как $\lambda_n \neq \lambda_m$, получаем:
\[ \int_0^l X_n(x)X_m(x) dx = 0. \]
\end{proof}

\textit{Замечание к Свойствам 1 и 4 (о счетности и асимптотике):} Доказательство счетности и того факта, что $\lambda_n \to +\infty$, строго вытекает из поиска нетривиальных решений вида $X(x) = A\cos(\sqrt{\lambda}x) + B\sin(\sqrt{\lambda}x)$. Подстановка этого вида в краевые условия приводит к трансцендентному уравнению на собственные значения (вида $\tan(\sqrt{\lambda}l) = \text{const}$). Графический анализ пересечения ветвей тангенса с прямыми или гиперболами дает счетное число корней, стремящихся к бесконечности. Строгое обоснование полноты системы $\{X_n(x)\}$ в пространстве $L_2(0,l)$ опирается на теорему Стеклова (из курса функционального анализа).

\subsection*{Обоснование выполнения краевых условий для ряда Фурье}

\begin{proof}
Покажем, что функция $v(x,t)$, сконструированная в виде суммы ряда $v(x,t) = \sum_{n=1}^\infty T_n(t) X_n(x)$, автоматически удовлетворяет однородным краевым условиям (3) и (4).
Предполагая, что ряд сходится достаточно хорошо и допускает почленное дифференцирование по пространственной переменной, подставим его в левое краевое условие ($x=0$):
\[ (\alpha v - \beta v_x)\big|_{x=0} = \alpha \sum_{n=1}^\infty T_n(t) X_n(0) - \beta \sum_{n=1}^\infty T_n(t) X_n'(0). \]
Внесем константы под знак суммы и сгруппируем слагаемые при $T_n(t)$:
\[ (\alpha v - \beta v_x)\big|_{x=0} = \sum_{n=1}^\infty T_n(t) \underbrace{\big[ \alpha X_n(0) - \beta X_n'(0) \big]}_{=0}. \]
Поскольку каждая функция $X_n(x)$ по построению является собственной функцией задачи Штурма-Лиувилля, выражение в квадратных скобках тождественно равно нулю для любого $n$. Следовательно, вся сумма обращается в нуль. 
Аналогичное рассуждение справедливо и для правого конца $x=l$:
\[ (\gamma v + \delta v_x)\big|_{x=l} = \sum_{n=1}^\infty T_n(t) \underbrace{\big[ \gamma X_n(l) + \delta X_n'(l) \big]}_{=0} = 0. \]
Таким образом, итоговый ряд является корректным решением с точки зрения выполнения нулевых граничных условий задачи.
\end{proof}

\subsection*{Обоснование Шага 3: Разделение переменных (вывод ОДУ для $T_n(t)$)}

\begin{proof}
Рассмотрим волновое уравнение $v_{tt} = a^2 v_{xx} + g(x,t)$ (для уравнения теплопроводности рассуждения аналогичны, с заменой $v_{tt}$ на $v_t$). Будем искать решение в виде суммы ряда:
\[ v(x,t) = \sum_{n=1}^\infty T_n(t) X_n(x). \]
Подставим этот ряд и разложение функции $g(x,t) = \sum_{n=1}^\infty c_n(t) X_n(x)$ в дифференциальное уравнение:
\[ \sum_{n=1}^\infty T_n''(t) X_n(x) = a^2 \sum_{n=1}^\infty T_n(t) X_n''(x) + \sum_{n=1}^\infty c_n(t) X_n(x). \]
Поскольку функции $X_n(x)$ являются решениями задачи Штурма-Лиувилля, для них выполняется тождество $X_n''(x) = -\lambda_n X_n(x)$. Сделаем замену:
\[ \sum_{n=1}^\infty T_n''(t) X_n(x) = - a^2 \sum_{n=1}^\infty \lambda_n T_n(t) X_n(x) + \sum_{n=1}^\infty c_n(t) X_n(x). \]
Перенесем все члены в левую часть и сгруппируем по базисным функциям $X_n(x)$:
\[ \sum_{n=1}^\infty \left[ T_n''(t) + a^2 \lambda_n T_n(t) - c_n(t) \right] X_n(x) = 0. \]
Так как система собственных функций $\{X_n(x)\}_{n=1}^\infty$ образует ортогональный базис, равенство ряда нулю тождественно для всех $x \in (0,l)$ возможно тогда и только тогда, когда коэффициенты при каждой функции $X_n(x)$ равны нулю. Отсюда получаем бесконечную систему независимых обыкновенных дифференциальных уравнений:
\[ T_n''(t) + a^2 \lambda_n T_n(t) = c_n(t), \quad \forall n \in \mathbb{N}. \]
Начальные условия для $T_n(t)$ выводятся аналогично из разложения условий $v(x,0) = v_0(x)$ и $v_t(x,0) = v_1(x)$:
\[ v(x,0) = \sum_{n=1}^\infty T_n(0) X_n(x) = \sum_{n=1}^\infty a_n X_n(x) \implies T_n(0) = a_n. \]
\[ v_t(x,0) = \sum_{n=1}^\infty T_n'(0) X_n(x) = \sum_{n=1}^\infty b_n X_n(x) \implies T_n'(0) = b_n. \]
Таким образом, поиск функции двух переменных $v(x,t)$ строго сведен к решению счетного набора задач Коши для обыкновенных дифференциальных уравнений.
\end{proof}

\end{document}